% =============================================================================
% 01_abstract.tex — Abstract / Zusammenfassung
% =============================================================================
% Prägnante Zusammenfassung der gesamten Arbeit (max. eine halbe A4-Seite).
% Inhalt: Was wurde analysiert? Welches Vorgehen wurde gewählt?
%         Welche Erkenntnisse wurden gewonnen?
% Der Abstract richtet sich an Lesende, die die Arbeit noch nicht kennen.
% =============================================================================

\chapter*{Abstract}
\addcontentsline{toc}{chapter}{Abstract}  % Manuell ins Inhaltsverzeichnis einfügen

% --- Analysiertes Problem / Ausgangslage ------------------------------------
% Beschreibe in 2–3 Sätzen, welches Problem oder welche Fragestellung
% den Ausgangspunkt der Arbeit bildet.
<<Hier die Ausgangslage und das analysierte Problem beschreiben.>>

% --- Gewähltes Vorgehen ------------------------------------------------------
% Erkläre kurz, wie das Problem angegangen wurde (Methoden, Technologien).
<<Hier das methodische Vorgehen und die eingesetzten Technologien beschreiben.>>

% --- Erzielte Erkenntnisse / Ergebnisse --------------------------------------
% Fasse die wichtigsten Resultate und Erkenntnisse der Arbeit zusammen.
<<Hier die wichtigsten Ergebnisse und Erkenntnisse zusammenfassen.>>
